\label{chp:rslts12}

This section is devoted on the analysis of the sum rules for the transition formfactors $G_1(0)$ and $G_2(0)$. For the values, the input parameters appearing in the sum rules are

\begin{table}[!ht]
    \centering
    \begin{tabular}{crl}
        $m_c\left(\overline{m}_c\right)$ &=& $\left(1.275 \pm 0.025\right) \text{GeV}$ \cite{PDG} \\
        $m_b\left(\overline{m}_b\right)$ &=& $\left(4.18 \pm 0.03\right) \text{GeV}$ \cite{PDG} \\
        $f_{3\gamma}$ &=& $-0.0039 \text{GeV}^2$ \cite{PhotonDA} \\
        $\chi$ &=& $\left(3.15\pm0.10\right) \text{GeV}^2$ \cite{PhotonDA} \\
        $\qbarq$ &=& $\left(-0.24 \pm 0.001\right)^3 \text{GeV}^3$ \cite{ioffeCondensate} \\
        $m_0^2$ &=& $\left(0.8 \pm 0.2\right) \text{GeV}^2$ \cite{SVZ}
    \end{tabular}
    \caption*{}
    \label{inputParams}
\end{table}
%
For the heavy quark masses, we have used their $\overline{\text{MS}}$ Scheme values. The masses of the spin-$\nicefrac{3}{2}$ doubly heavy baryons and spin-$\nicefrac{1}{2}$ doubly heavy baryons are calculated in \cite{lattice_masses,AlievOctet,AlievDecuplet}, and the mass of $\Xi_{cc}$ baryon is experimentally observed.

\begin{table}[ht]
    \centering
    \begin{tabular}{|c|c|c|c|c|}
         \hline
         Baryon & Lattice\cite{lattice_masses} &  QCDSR\cite{AlievOctet} & QCDSR\cite{AlievDecuplet} & Exp\cite{xi+,xi++}\\
         \hline
         $\Xi_{bc}$ & 6.943 GeV & $6.72 GeV$ & - & -  \\
         \hline
         $\Xi_{bc}^{\prime}$ & 6.959 GeV & $6.79 GeV$ & - & - \\
         \hline
         $\Xi_{bb}$ & 10.143 GeV & $9.96 GeV$ & - & - \\
         \hline
         $\Xi_{cc}$ & 3.610 GeV& $3.72 GeV$ & - & $3.52 GeV$  \\
         \hline
         $\Xi^{\ast}_{bc}$ & 6.985 GeV& - & $7.25 GeV$& -  \\
         \hline
         $\Xi^{\ast}_{cc}$ & 3.692 GeV& - &  $3.69  GeV$& - \\
         \hline
         $\Xi^{\ast}_{bb}$ & 10.178 GeV & - & $10.4  GeV$& -  \\
         \hline
    \end{tabular}
    \caption{Baryon masses}
    \label{masstable}
\end{table}
We will see that the radiative decay widths of spin-$\nicefrac{3}{2}$ to spin-$\nicefrac{1}{2}$ baryons are proportional to the cube of the mass difference of the initial and final baryons, $\left(\Delta m\right) = \left(m_{\nicefrac{3}{2}}-m_{\nicefrac{1}{2}}\right)$. Therefore, the decay widths are very sensitive to the mass difference of doubly heavy baryons. For this reason, in the next discussion, we will use the mass of doubly heavy baryons obtained from lattice calculations which contain a minimal error.
The photon distribution amplitudes are the main non-perturbative inputs of light cone sum rules. The expression of DAs and the values of the parameters entering in the expressions of DAs are given in \cite{PhotonDA}. Now, let us perform the numerical analysis of the relevant formfactors.

The sum rules for the transition formfactors $G_1(0)$ and $G_2(0)$ involve three auxiliary parameters: parameter $\beta$ appearing in the expression of the interpolating current for spin-$\nicefrac{1}{2}$ doubly heavy baryons, the continuum threshold $s_0$ and Borel mass parameter $M^2$. The working region of $s_0$ is determined from the analysis of two-point sum rules which are carried out in \cite{AlievOctet,AlievDecuplet} and given in Table-\ref{parameter_table}

The working region of $M^2$ is determined from two requirements. From one side $M^2$ must be large enough to guarantee the dominance of leading twist over higher twist contributions, and from the other side, it should be small enough in order to ensure the suppression of the higher states and continuum contributions. The working regions of Borel mass parameters are also presented in Table-\ref{parameter_table}.

As an example, in Figures-\ref{xibcpgms}(a) and \ref{xibcpges}(a) we present the dependencies of $G_M(0)$ and $G_E(0)$ on $M^2$ at fixed values of $s_0$ at $\beta$. From these figures, we see that the formfactors $G_M$ and $G_E$ exhibit excellent stability with respect to the variation of $M^2$ in its working region.

\tikzexternalenable
\label{chp:appndxd}

\pgfplotsset{width=10cm,compat=1.18}


%-------------------------------------------------------------------------------------------------% 

\begin{figure}[h!]
    \centering
    \subfigure[]{
    \begin{tikzpicture}[scale=0.6]
		\begin{axis}[
			xlabel={$M^2 \left(GeV^2\right)$},
			ylabel={$|G_M| \left(\mu_N\right)$},
			enlargelimits=false,
			ymin=0.6,
			ymax=2,
			]
			\addplot [
			color=black,
			mark=square,
			mark size=2.9pt]
			table[x={Msq}, y={gmlow}]{dataforfigs/xibcpgs.dat};
			\addlegendentry{$s_0 = 56 GeV^2$}
			\addplot [
			color=black,
			mark=o,
			mark size=2.9pt]
			table[x={Msq}, y={gmmid}]{dataforfigs/xibcpgs.dat};
			\addlegendentry{$s_0 = 57 GeV^2$}
			\addplot [
			color=black,
			mark=triangle,
			mark size=2.9pt]
			table[x={Msq}, y={gmhigh}]{dataforfigs/xibcpgs.dat};
			\addlegendentry{$s_0 = 58 GeV^2$}
		\end{axis}
	\end{tikzpicture}
    }
    \subfigure[]{
    \begin{tikzpicture}[scale=0.6]
    \pgfplotsset{
            every axis legend/.append style={
                at={(0.3,0.025)},
                anchor=south west,
            },
        }
		\begin{axis}[
			xlabel={$\cos\left(\theta\right)$},
			ylabel={$|G_M| \left(\mu_N\right)$},
			enlargelimits=false,
			ymin=0,
			ymax=2,
			]
			\addplot [
			color=black,
			mark=square,
			mark repeat=10,
			mark phase=10,
			mark size=2.9pt]
			table[x={cos}, y={gm58}]{dataforfigs/xibcpcosthetas.dat};
			\addlegendentry{$s_0 = 58 GeV^2$}
			\addplot [
			color=black,
			mark=o,
			mark repeat=10,
			mark phase=10,
			mark size=2.9pt]
			table[x={cos}, y={gm59}]{dataforfigs/xibcpcosthetas.dat};
			\addlegendentry{$s_0 = 59 GeV^2$}
			\addplot [
			color=black,
			mark=triangle,
			mark repeat=10,
			mark phase=10,
			mark size=2.9pt]
			table[x={cos}, y={gm60}]{dataforfigs/xibcpcosthetas.dat};
			\addlegendentry{$s_0 = 60 GeV^2$}
		\end{axis}
	\end{tikzpicture}
    }
    \caption[$G_M$ for $\Xi^{\ast}+_{bc}\rightarrow\Xi^{+}_{bc}\gamma$]{\textbf{a)} $M^2$ dependence of the $G_M$ for $\Xi^{\ast+}_{bc}\rightarrow\Xi^{+}_{bc}\gamma$ decay with different values of $s_0$ at $\beta=5$, \textbf{b)} $\cos\left(\theta\right)$ dependence of $G_M$ for $\Xi^{\ast+}_{bc}\rightarrow\Xi^{+}_{bc}\gamma$ with different values of $s_0$ at $M^2=8$ $GeV^2$ where $\beta=\tan\left(\theta\right)$}
    \label{xibcpgms}
\end{figure}

\begin{figure}[h!]
    \centering
    \subfigure[]{
    \begin{tikzpicture}[scale=0.6]
		\begin{axis}[
			xlabel={$M^2 \left(GeV^2\right)$},
			ylabel={$|G_E| \left(10^{-3}\mu_N\right)$},
			enlargelimits=false,
			ymin=0,
			ymax=5,
			]
			\addplot [
			color=black,
			mark=square,
			mark size=2.9pt]
			table[x={Msq}, y={gelow}]{dataforfigs/xibcpges.dat};
			\addlegendentry{$s_0 = 58 GeV^2$}
			\addplot [
			color=black,
			mark=o,
			mark size=2.9pt]
			table[x={Msq}, y={gemid}]{dataforfigs/xibcpges.dat};
			\addlegendentry{$s_0 = 59 GeV^2$}
			\addplot [
			color=black,
			mark=triangle,
			mark size=2.9pt]
			table[x={Msq}, y={gehigh}]{dataforfigs/xibcpges.dat};
			\addlegendentry{$s_0 = 60 GeV^2$}
		\end{axis}
	\end{tikzpicture}  
    }
	\subfigure[]{
    \begin{tikzpicture}[scale=0.6]
		\begin{axis}[
			xlabel={$\cos\left(\theta\right)$},
			ylabel={$|G_E| \left(10^3\mu_N\right)$},
			enlargelimits=false,
			ymin=0,
			ymax=5,
			]
			\addplot [
			color=black,
			mark=square,
			mark repeat=10,
			mark phase=10,
			mark size=2.9pt]
			table[x={cos}, y={ge58}]{dataforfigs/xibcpcosthetas.dat};
			\addlegendentry{$s_0 = 58 GeV^2$}
			\addplot [
			color=black,
			mark=o,
			mark repeat=10,
			mark phase=10,
			mark size=2.9pt]
			table[x={cos}, y={ge59}]{dataforfigs/xibcpcosthetas.dat};
			\addlegendentry{$s_0 = 59 GeV^2$}
			\addplot [
			color=black,
			mark=triangle,
			mark repeat=10,
			mark phase=10,
			mark size=2.9pt]
			table[x={cos}, y={ge60}]{dataforfigs/xibcpcosthetas.dat};
			\addlegendentry{$s_0 = 60 GeV^2$}
		\end{axis}
	\end{tikzpicture}
    }
    \caption[$G_E$ for $\Xi^{\ast+}_{bc}\rightarrow\Xi^{+}_{bc}\gamma$]{Same as \ref{xibcpgms} but for $G_E$}
    \label{xibcpges}
\end{figure}

Finally in order to find the working region of $\beta$, we present the dependence of $G_M(0)$ and $G_E(0)$ for the $\Xi_{bc}^{\ast+}\rightarrow\Xi_{bc}^{+}$ transition on $\cos(\theta)$ where $\beta=\tan(\theta)$ in Figures-\ref{xibcpgms}(b) and \ref{xibcpges}(b). From these figures, we see that the results on the formfactors show very good stability when $-0.4<\cos(\theta)<0.4$. Performing similar calculations for the formfactors of other transitions we find that in this region of $\cos(\theta)$ the formfactors are rather stable. Therefore the working region of $\cos\theta$ is $\left(-0.4,0.4\right)$

\begin{table}[ht]
    \centering
    \begin{tabular}{|c|c|c|c|}
        \hline
         Transition & $M^2$ & $\cos(\theta)$ & $s_0$  \\
         \hline
         $\Xi_{bc}^+\rightarrow\Xi_{bc}^{\ast+}$ & $6-9~GeV^2$ & $(-0.4,0.4)$ & $57\pm1 ~GeV^2$\\
         \hline
         $\Xi_{bc}^{\prime+}\rightarrow\Xi_{bc}^{\ast+}$& $6-9~GeV^2$ & $(-0.4,0.4)$ & $57\pm1~ GeV^2$\\
         \hline
         $\Xi_{bc}^0\rightarrow\Xi_{bc}^{\ast0}$& $6-9~GeV^2$ & $(-0.4,0.4)$ & $57\pm1~ GeV^2$\\
         \hline
         $\Xi_{bc}^{\prime0}\rightarrow\Xi_{bc}^{\ast0}$& $6-9 ~GeV^2$ & $(-0.4,0.4)$ & $57\pm1~ GeV^2$\\
         \hline
         $\Xi_{cc}^{++}\rightarrow\Xi_{cc}^{\ast++}$& $3-6 ~GeV^2$ & $(-0.4,0.4)$ & $20\pm1~ GeV^2$\\
         \hline
         $\Xi_{cc}^+\rightarrow\Xi_{cc}^{\ast+}$& $3-6 ~GeV^2$ & $(-0.4,0.4)$ & $20\pm1 ~GeV^2$\\
         \hline
         $\Xi_{bb}^0\rightarrow\Xi_{bb}^{\ast0}$& $9-12 ~GeV^2$ & $(-0.4,0.4)$ & $121\pm2 ~GeV^2$\\
         \hline
         $\Xi_{bb}^-\rightarrow\Xi_{bb}^{\ast-}$& $9-12 ~GeV^2$ & $(-0.4,0.4)$ & $121\pm2 ~GeV^2$\\
        \hline
    \end{tabular}
    \caption{Working Regions for $M^2$, $\beta$ and $s_0$}
    \label{parameter_table}
\end{table}
We performed the numerical analysis by using the working regions of $\cos\theta$, $M^2$ and $s_0$, and our final findings of the moments $G_M(0)$ and $G_E(0)$ are collected in Table-\ref{magnetic_results} and Table-\ref{electric_results}. Note that, since the signs of the residues are undetermined, we present the absolute values of our results.
\begin{table}[!ht]
    \centering
    \begin{tabular}{|c|c|c|c|c|c|c|c|}
        \hline
         Transition & This Work & QM\cite{transition15} & $\chi$ PT\cite{transition15} & Bag Model\cite{transition10} & EMS\cite{DhirEMS} & exBag\cite{exBag} & Lattice\cite{lattice} \\
         \hline
         $\Xi_{bc}^+\rightarrow\Xi_{bc}^{\ast+}$& $1.25\pm0.156$ & $-1.61$ & -2.56 & 0.695 &  1.12 & 1.12 & - \\
         \hline
         $\Xi_{bc}^{\prime+}\rightarrow\Xi_{bc}^{\ast+}$& $0.17\pm0.030$ & $-0.36$ & -0.36 & 0.672  & 1.04 & 0.814 & - \\
         \hline
         $\Xi_{bc}^0\rightarrow\Xi_{bc}^{\ast0}$& $0.77\pm0.104$ & $1.02$ &  1.03 & -0.747 & -1.03 & -0.919 & - \\
         \hline
         $\Xi_{bc}^{\prime0}\rightarrow\Xi_{bc}^{\ast0}$& $0.18\pm0.032$ & $-0.36$  & -0.36 & 0.070 & -0.17 & 0.598 & - \\
         \hline
         $\Xi_{cc}^{++}\rightarrow\Xi_{cc}^{\ast++}$& $1.03\pm0.138$ & $-1.4$  & -2.35 & -0.787 & -1.30 & -1.21 & -0.772\\
         \hline
         $\Xi_{cc}^+\rightarrow\Xi_{cc}^{\ast+}$& $0.96\pm0.158$ & $1.23$  & 1.55 & 0.945 & 1.19 & 1.07 & 0.906\\
         \hline
         $\Xi_{bb}^0\rightarrow\Xi_{bb}^{\ast0}$& $1.78\pm0.300$ & $-1.82$ & -2.77 & -1.039 & -1.70 & -1.045 & - \\
         \hline
         $\Xi_{bb}^-\rightarrow\Xi_{bb}^{\ast-}$& $0.82\pm0.131$ & $0.81$ & 1.13 & 0.428 & 0.76 & 0.643 & - \\
        \hline
    \end{tabular}
    \caption{Transition magnetic moments in nuclear magneton}
    \label{magnetic_results}
\end{table}
\begin{table}[ht!]
    \centering
    \begin{tabular}{|c|c|}
        \hline
         Transition & $|G_E|$  \\
         \hline
         $\Xi_{bc}^+\rightarrow\Xi_{bc}^{\ast+}$& $\left( 1.92\pm0.239\right)\times10^{-3}$\\
         \hline
         $\Xi_{bc}^{\prime+}\rightarrow\Xi_{bc}^{\ast+}$& $\left( 0.167\pm0.029\right)\times10^{-3}$\\
         \hline
         $\Xi_{bc}^0\rightarrow\Xi_{bc}^{\ast0}$& $\left( 1.18\pm0.160\right)\times10^{-3}$\\
         \hline
         $\Xi_{bc}^{\prime0}\rightarrow\Xi_{bc}^{\ast0}$& $\left( 0.18\pm0.032\right)\times10^{-3}$\\
         \hline
         $\Xi_{cc}^{++}\rightarrow\Xi_{cc}^{\ast++}$& $\left( 5.81\pm0.775\right)\times10^{-3}$\\
         \hline
         $\Xi_{cc}^+\rightarrow\Xi_{cc}^{\ast+}$& $\left( 5.44\pm0.901\right)\times10^{-3}$\\
         \hline
         $\Xi_{bb}^0\rightarrow\Xi_{bb}^{\ast0}$& $\left( 1.56\pm0.264\right)\times10^{-3}$ \\
         \hline
         $\Xi_{bb}^-\rightarrow\Xi_{bb}^{\ast-}$& $\left( 0.72\pm0.115\right)\times10^{-3}$ \\
        \hline
    \end{tabular}
    \caption{Transition electric quadrupole moments in nuclear magneton}
    \label{electric_results}
\end{table}
 All errors coming from uncertainties of input parameters, as well as the variation of $s_0$, $M^2$, and $\cos\theta$ in their working regions, are taken quadratically. For comparison in Table-\ref{magnetic_results}, we present the predictions for the magnetic moments $G_M(0)$ of other approaches such as quark model \cite{transition15}, $\chi$PT \cite{transition15}, bag model \cite{transition10}, effective quark mass scheme \cite{DhirEMS}, extended bag model \cite{exBag} and lattice QCD \cite{lattice}
Using the results on the moments $G_M$ and $G_E$ one can easily estimate the decay widths with the help of the following formula
\begin{equation}
    \Gamma = \dfrac{\alpha}{16}\dfrac{\left(m_1^2-m_2^2\right)^3}{m_N^2m_1^3}\left[G_M^2(0)+3G_E^2(0)\right]
    \label{decay_width_formula}
\end{equation}
where $m_N$ is the mass of the nucleon.
The results of the decay widths are collected in table-\ref{decays_table}.
\begin{table}[ht]
    \centering
    \begin{tabular}{|c|c|c|c|c|c|c|}
        \hline
         Decay & This Work & $\chi$ PT\cite{transition15} & Bag Model\cite{transition10} & exBag Model\cite{exBag} & RQM\cite{RQM} & Lattice\cite{lattice}  \\
         \hline
         $\Xi_{bc}^{\ast+}\rightarrow\Xi_{bc}^{+}\gamma$& $0.48\pm0.119$ & $26.2$ & 0.533 & 1.31 & - & -\\
         \hline
         $\Xi_{bc}^{\ast+}\rightarrow\Xi_{bc}^{\prime+}\gamma$& $0.002\pm0.0007$ & $0.52$ & 0.031 & 0.0293 & - & -\\
         \hline
         $\Xi_{bc}^{\ast0}\rightarrow\Xi_{bc}^{0}\gamma$& $0.18\pm0.049$ & $7.19$ &  0.612 & 0.876 & - & -\\
         \hline
         $\Xi_{bc}^{\ast0}\rightarrow\Xi_{bc}^{\prime0}\gamma$& $0.003\pm0.0009$ & $0.52$  & 0.000 & $7.6 \cdot 10^{-5}$ & - & -\\
         \hline
          $\Xi_{cc}^{\ast++}\rightarrow\Xi_{cc}^{++}\gamma$& $2.36\pm0.622$ & 22.0 & 1.43 & 2.79 & 7.21 & 0.0518\\
         \hline
         $\Xi_{cc}^{\ast+}\rightarrow\Xi_{cc}^{+}\gamma$& $2.07\pm0.666$ & 9.57 & 2.08 & 2.17 & 3.90 & 0.0648\\
         \hline
         $\Xi_{bb}^{\ast0}\rightarrow\Xi_{bb}^{0}\gamma$& $0.58\pm0.188$ & $31.1$ & 0.126 & 0.137 & 0.98 & -\\
         \hline
         $\Xi_{bb}^{\ast-}\rightarrow\Xi_{bb}^{-}\gamma$& $0.12\pm0.038$ & $5.17$ & 0.022 & 0.0268 & 0.21 & -\\
        \hline
    \end{tabular}
    \caption{Decay widths in keV}
    \label{decays_table}
\end{table}

For completeness, we also presented the decay width results obtained from the bag model,exBag model, RQM, and $\chi$PT. We see that our results are drastically different from the results of both approaches except the results for $\Xi_{bc}^{\ast+}\longrightarrow\Xi_{bc}^+\gamma$ and $\Xi_{bc}^{\ast0}\longrightarrow\Xi_{bc}^0\gamma$ transition obtained from bag model which are in agreement with our results. 

From Table-\ref{magnetic_results} we get the following results: Our finding of transition magnetic moments for $\Xi_{bc}^{\ast+}\rightarrow\Xi_{bc}^{+}$, $\Xi_{cc}^{++}\rightarrow\Xi_{cc}^{\ast++}$, $\Xi_{cc}^{+}\rightarrow\Xi_{cc}^{\ast+}$ and $\Xi_{bb}^{\ast-}\rightarrow\Xi_{bb}^{-}$ transitions are in agreement with the QM\cite{transition15}, EMS \cite{DhirEMS} and exBag \cite{exBag} but considerably differ from $\chi$PT\cite{transition15}. Results for the $\Xi_{bc}^{\ast0}\rightarrow\Xi_{bc}^{0}$ transition is in agreement with $\chi$PT \cite{transition15} and exBag \cite{exBag}, $\Xi_{bc}^{\ast0}\rightarrow\Xi_{bc}^{\prime0}$ transition is agreeing with the EMS \cite{DhirEMS}, $\Xi_{bc}^{\ast0}\rightarrow\Xi_{bc}^{0}$ transition is in agreement with bag model \cite{transition10} and exBag \cite{exBag}, and  $\Xi_{bb}^{\ast0}\rightarrow\Xi_{bb}^{0}$ agrees with the both QM \cite{transition15} and EMS \cite{DhirEMS}. While the results for $\Xi_{cc}^{++}\rightarrow\Xi_{cc}^{\ast++}$ also agrees with the quark model \cite{transition15} and $\Xi_{cc}^{+}\rightarrow\Xi_{cc}^{\ast+}$ agrees with the lattice qcd \cite{lattice}, results for $\Xi_{bc}^{\ast+}\rightarrow\Xi_{bc}^{\prime+}$ transition differs from other approaches results.

From Table-\ref{decays_table} we get the following results: Our finding of the decay widths for the $\Xi_{bc}^{\ast+}\rightarrow\Xi_{bc}^{+}$ transition is in agreement with the bag model\cite{transition10}, $\Xi_{cc}^{\ast++}\rightarrow\Xi_{cc}^{++}$ and $\Xi_{cc}^{\ast+}\rightarrow\Xi_{cc}^{+}$ are in agreement with the exBag model\cite{exBag} while $\Xi_{cc}^{\ast+}\rightarrow\Xi_{cc}^{+}$ transition also agrees with the bag model\cite{transition10}. For other transitions, the results differ from other works. As we stated earlier, the values of the decay widths are sensitive to the mass difference of the initial and final baryons, hence the results differ drastically with the selection of the masses.



\section{Conclusion}

We calculated the magnetic dipole $G_M(0)$ and electric quadrupole $G_E(0)$ formfactors of spin-$\nicefrac{3}{2}$ doubly heavy baryons to spin-$\nicefrac{1}{2}$ doubly heavy baryon transitions within light cone sum rules method. Having the values of the magnetic dipole $G_M(q^2)$ and electric quadrupole $G_E(q^2)$ formfactors at $q^2=0$ we estimate the corresponding decay widths. Our findings on formfactors and decay widths compared with the predictions of other approaches.   